\label{ch:results}

This chapter presents, interprets, and summarises the real numerical
results of \texttt{inference6\_sno\_ii.ipynb} (in
\texttt{notebooks/inference/}), executed end to end against the real
391-day salt-phase 38-observable data (\cref{ch:data-release}) using the
physics of \cref{ch:detector-implementation} and the inference machinery
of \cref{ch:inference-process}. Every number below is a real, reproducible
output of that notebook, not illustrative.

\section{Real data and covariance}

\begin{figure}[htb]
\centering
\includegraphics[width=0.85\textwidth]{fig1_cc_spectrum}
\caption{Real published day/night CC equivalent-flux spectrum, Table XXX, statistical error only.}
\label{fig:cc-spectrum}
\end{figure}

\Cref{fig:cc-spectrum} shows the real 17-bin CC spectrum loaded directly
from \texttt{data/detector/sno\_ii/}. \Cref{tab:integrated-fluxes} lists
the real integrated NC and ES fluxes with their published statistical and
systematic uncertainties separately.

\begin{table}[htb]
\centering
\caption{Real integrated equivalent fluxes (Table XXIV), units $10^6\,{\rm cm^{-2}s^{-1}}$.}
\label{tab:integrated-fluxes}
\begin{tabular}{@{}lrrrr@{}}
\toprule
Channel & Day & $\sigma_{\rm stat}$ & $\sigma_{\rm syst}$ & Night ($\sigma_{\rm stat}$, $\sigma_{\rm syst}$) \\
\midrule
NC & 4.81 & 0.31 & 0.39 & 5.02 (0.29, 0.41) \\
ES & 2.17 & 0.14 & 0.14 & 2.52 (0.32, 0.16) \\
\bottomrule
\end{tabular}
\end{table}

\begin{resultbox}
The 17 real day (night) CC bins sum to $1.7325$ ($1.6422$), matching the
independently published integral flux, $1.73$ ($1.64$) -- confirming the
CC spectral table and the integrated-flux table are mutually consistent,
as SNO's own Table XXX caption states they must be.
\end{resultbox}

\begin{figure}[htb]
\centering
\includegraphics[width=0.95\textwidth]{fig2_covariance_matrices}
\caption{Correlation matrices for $V_{\rm stat}$ (left) and $V_{\rm stat}+V_{\rm syst}$ (right); black lines mark the day/night block boundary.}
\label{fig:covariance-matrices}
\end{figure}

\Cref{fig:covariance-matrices} shows both real correlation structures
(\cref{sec:covariance-construction}): $V_{\rm stat}$'s clear day/night
block-diagonal structure (statistically independent periods) contrasts
visibly with $V_{\rm stat}+V_{\rm syst}$'s non-zero off-block correlation,
the direct visual signature of shared detector systematics affecting both
periods identically.

\begin{resultbox}
Both covariance matrices are verified positive-definite before every fit
below: $V_{\rm stat}$'s eigenvalues span $[6.25\times10^{-6},\,
1.044\times10^{-1}]$; $V_{\rm stat}+V_{\rm syst}$'s span $[6.37\times10^{-6},\,
3.186\times10^{-1}]$ -- both comfortably positive, so
\texttt{torch.linalg.cholesky} never fails.
\end{resultbox}

\section{Detector-level closure checks}
\label{sec:results-closure}

\begin{resultbox}
Before any fit, the real forward model (built at NuFIT~6.1 defaults,
$\theta_{12}=33.76^\circ$, $\Delta m^2_{21}=7.537\times10^{-5}\,{\rm
eV}^2$, $\lambda_{8B}=0$) is checked against three real physics identities:
\begin{itemize}
  \item At $P_{ee}\equiv1$, the CC spectrum's own sum equals $\Phi_{\rm
    tot}$ to a relative difference of $0.0000\%$ (both quantities
    computed on the identical finite quadrature grid, so this is an exact
    self-consistency check of \labelcref{eq:cc-equivalent-flux}'s own
    ratio construction, not a claim about the grid's absolute accuracy).
  \item NC's real day/night values at NuFIT defaults, $4.1248$ vs.\
    $4.1249$, differ by $0.0005\%$ -- confirming
    \labelcref{eq:nc-equivalent-flux}'s Standard-Model oscillation
    independence numerically.
  \item Autograd gradients of the full prediction with respect to every
    one of the three free parameters are finite and non-zero,
    $\nabla=(10.47,\,11672.17,\,14.57)$ -- a real LBFGS fit is possible.
\end{itemize}
At these (non-optimized) NuFIT-default parameters, the forward model gives
NC day/night $=4.12/4.12$ (published: $4.81/5.02$) and ES day/night
$=1.78/1.82$ (published: $2.17/2.52$) -- the right order of magnitude, as
expected of a starting point rather than a fit.
\end{resultbox}

\section{Real fit: statistical covariance only}
\label{sec:results-fit-stat}

\begin{resultbox}
The first real fit against all 38 published observables, using
$V_{\rm stat}$ alone (\labelcref{eq:stat-covariance}), converges in 29
closure evaluations ($\approx$50~s):
\begin{equation}
  \keyresult{\theta_{12} = 35.69^\circ \pm 1.41^\circ, \qquad
    \Delta m^2_{21} = (5.684\pm0.945)\times10^{-5}\,{\rm eV}^2, \qquad
    \lambda_{8B} = 0.162\pm0.038,}
  \label{eq:result-vstat-fit}
\end{equation}
i.e.\ $\Phi_{8B}=4.862\times10^6\,{\rm cm^{-2}s^{-1}}$; $-2\ln L$ improves
from $154.94$ to $46.65$. $\theta_{12}$ sits comfortably inside SNO's own
published SNO-only $1\sigma$ range ($\approx30.6^\circ$--$36.8^\circ$,
Table XXVIII); $\Delta m^2_{21}$ sits just above SNO's own tight SNO-only
upper bound ($5.0^{+0.2}_{-1.8}\times10^{-5}$, i.e.\ $\le5.2\times10^{-5}$)
but well within its broader Global-solar range
($6.5^{+4.4}_{-2.3}\times10^{-5}$) -- \cref{sec:results-comparison}
compares this quantitatively and geometrically against the paper's own
published contour.
\end{resultbox}

\section{Real fit: statistical plus systematic covariance}
\label{sec:results-fit-total}

\begin{resultbox}
Adding $V_{\rm syst}$ (\labelcref{eq:syst-covariance}), the fit converges
to
\begin{equation}
  \theta_{12} = 35.28^\circ \pm 1.61^\circ, \qquad
  \Delta m^2_{21} = (6.095\pm349{,}000)\times10^{-5}\,{\rm eV}^2, \qquad
  \lambda_{8B} = 0.129\pm0.064,
  \label{eq:result-vtotal-fit}
\end{equation}
i.e.\ $\Phi_{8B}=4.705\times10^6\,{\rm cm^{-2}s^{-1}}$; $-2\ln L$ improves
from $50.31$ to $41.21$. Central values shift only modestly relative to
\labelcref{eq:result-vstat-fit} ($\theta_{12}$ by $-0.41^\circ$), as
expected from a positive-semi-definite addition to the covariance.
\end{resultbox}

\begin{notebox}
\textbf{A genuine methodological finding, stated plainly rather than
smoothed over.} The \emph{reported} $\Delta m^2_{21}$ uncertainty in
\labelcref{eq:result-vtotal-fit} explodes to $\pm0.349\,{\rm eV}^2$ --
around 5700 times the central value. This is \emph{not} a claim that
$V_{\rm syst}$ makes $\Delta m^2_{21}$ literally unconstrained: it is the
Laplace/Fisher approximation (\labelcref{eq:laplace-covariance}-style
construction, a \emph{local}, Gaussian uncertainty estimate) breaking down
where the $-2\ln L$ surface has gone nearly flat in that direction once
$V_{\rm syst}$'s large, highly-correlated (rank $\le16$ of 38,
\cref{sec:covariance-construction}) contribution is added --
\pyfunc{fit\_lbfgs}'s own documented eigenvalue floor (``an unconstrained
direction gets a large but finite sigma instead of \code{nan}'') is doing
exactly its documented job here, not failing silently. This is precisely
why \cref{sec:results-scan}'s profiled likelihood scan, not this section's
single Laplace $\sigma$, is the statistically appropriate tool for
$\Delta m^2_{21}$'s real uncertainty once systematics are included.
\end{notebox}

\section{2-D likelihood scan}
\label{sec:results-scan}

Two variants of the same $(\theta_{12},\Delta m^2_{21})$ map, both
centred on \labelcref{eq:result-vtotal-fit}'s own best fit, using
$V_{\rm stat}+V_{\rm syst}$ throughout: \textbf{fixed-slice} holds
$\lambda_{8B}$ at its own best-fit value at every grid point (one forward
pass each); \textbf{profiled} re-minimizes $\lambda_{8B}$ at every grid
point (one LBFGS run each, $\approx$16-20~s/point).

\begin{figure}[htb]
\centering
\includegraphics[width=0.72\textwidth]{fig3_fixedslice_scan}
\caption{Fixed-slice $-2\ln L$ contours, $15\times15$ grid, $\lambda_{8B}$ fixed at its own best fit.}
\label{fig:fixedslice-scan}
\end{figure}

\begin{figure}[htb]
\centering
\includegraphics[width=0.72\textwidth]{fig4_profiled_scan}
\caption{Profiled $-2\ln L$ contours, deliberately coarse $5\times5$ grid, $\lambda_{8B}$ re-minimized per point.}
\label{fig:profiled-scan}
\end{figure}

\begin{resultbox}
\Cref{fig:fixedslice-scan} (164~s total, $15\times15$ points) and
\cref{fig:profiled-scan} (260~s total, deliberately coarse $5\times5$
points -- computational cost, not a statistical choice) both place their
own minimum at, or immediately adjacent to, \labelcref{eq:result-vtotal-fit}'s
best fit, as required for internal consistency. The qualitative shape --
elongated toward larger $\theta_{12}$ at larger $\Delta m^2_{21}$ -- matches
SNO's own published SNO-only contour shape (their Figure~34)
qualitatively (\cref{sec:results-comparison}).
\end{resultbox}

\section{Day-night sensitivity}

\begin{figure}[htb]
\centering
\includegraphics[width=0.78\textwidth]{fig5_daynight_asymmetry}
\caption{Predicted CC day-night asymmetry $A_{\rm CC}$ vs.\ $\Delta m^2_{21}$, at the fitted $\theta_{12}$.}
\label{fig:daynight-asymmetry}
\end{figure}

\Cref{fig:daynight-asymmetry} scans $A_{\rm CC}=2(\Phi_N-\Phi_D)/(\Phi_N+\Phi_D)$
across $\Delta m^2_{21}$ at the fitted $\theta_{12}$, confirming the
expected qualitative trend: $A_{\rm CC}$ grows with $\Delta m^2_{21}$ (more
Earth-matter regeneration at larger mass-splitting within the range
probed) and vanishes at very small $\Delta m^2_{21}$. SNO's own real
measured asymmetry, $A_{\rm CC}=-0.021\pm0.063({\rm stat})\pm0.035({\rm
syst})$ (Eq.~17 of \cite{Aharmim2005SaltPhase}), is statistics-limited and
consistent with zero; this scan illustrates the predicted \emph{trend}, not
a claim that SNO alone resolves a non-zero asymmetry.

\section{Comparison to SNO's own published best fit and contours}
\label{sec:results-comparison}

\begin{theorybox}
SNO quotes $\tan^2\theta$, not $\theta_{12}$ directly (Table XXVIII of
\cite{Aharmim2005SaltPhase}); the conversion
$\theta_{12}=\arctan\sqrt{\tan^2\theta}$ is applied here, including to the
asymmetric $1\sigma$ bounds, propagated through the same nonlinear
transform rather than treated as symmetric. \Cref{tab:table-xxviii}
lists the three published rows converted this way.
\end{theorybox}

\begin{table}[htb]
\centering
\caption{Table XXVIII of \cite{Aharmim2005SaltPhase}, converted to $\theta_{12}$.}
\label{tab:table-xxviii}
\begin{tabular}{@{}lrr@{}}
\toprule
Analysis & $\theta_{12}$ (deg) & $\Delta m^2_{21}$ ($10^{-5}$ eV$^2$) \\
\midrule
SNO-only (LMA region only) & $33.85^{+2.95}_{-3.25}$ & $5.0^{+0.2}_{-1.8}$ \\
Global solar & $33.85^{+2.46}_{-2.54}$ & $6.5^{+4.4}_{-2.3}$ \\
Solar + KamLAND & $33.85^{+2.46}_{-2.20}$ & $8.0^{+0.6}_{-0.4}$ \\
\midrule
\textbf{This document (\cref{sec:results-fit-total})} & \textbf{35.28} & \textbf{6.10} \\
\bottomrule
\end{tabular}
\end{table}

\begin{figure}[htb]
\centering
\includegraphics[width=0.82\textwidth]{fig6_comparison_published}
\caption{This document's own fixed-slice contour (\cref{fig:fixedslice-scan}) with SNO's own three published best fits and $1\sigma$ boxes overlaid.}
\label{fig:comparison-published}
\end{figure}

\begin{resultbox}
\Cref{fig:comparison-published} overlays this document's own contour with
SNO's three published points. This document's fitted point falls
\emph{within} SNO's own SNO-only $1\sigma$ box in $\theta_{12}$, but
\emph{not} in $\Delta m^2_{21}$ (just above its tight upper bound) --
stated plainly rather than smoothed over. The three published rows
themselves sit at essentially the same $\theta_{12}\approx33.9^\circ$
(they share $\tan^2\theta=0.45$ almost exactly) but at increasingly
precise, increasingly high $\Delta m^2_{21}$ as more data is combined in
-- the well-known real physics story of $\Delta m^2_{21}$ being
\emph{KamLAND's} discovery, not a salt-phase-SNO-alone one, visible
directly in this figure's three published markers.
\end{resultbox}

\begin{notebox}
\textbf{On the comparison method.} SNO's own contour shape (their
Figure~34) is available here only as a rendered PDF image, not
machine-readable digitized points; overlaying a precisely digitized
contour line would overstate this document's own precision in reading it.
\Cref{tab:table-xxviii,fig:comparison-published} instead use the more
defensible comparison available: SNO's own \emph{numerical} best fit and
asymmetric $1\sigma$ box (exact, not read off a plot).
\end{notebox}

\section{Summary and known limitations}

This document fits SNO's real published salt-phase 391-day data set using
its own officially extracted 38-observable vector and full correlated
statistical+systematic covariance (Eq.~19-21 of
\cite{Aharmim2005SaltPhase}), reproducing the collaboration's own
likelihood construction rather than approximating it with a diagonal one.
Two real fits (\cref{sec:results-fit-stat,sec:results-fit-total}), both
converging to physically sensible $(\theta_{12},\Delta m^2_{21},\Phi_{8B})$
points, are compared quantitatively and geometrically against SNO's own
published best fit (\cref{sec:results-comparison}); a genuine
methodological finding not anticipated going in -- the Laplace
approximation's breakdown once $V_{\rm syst}$ is added
(\cref{sec:results-fit-total}) -- directly motivates and validates the
fixed-slice/profiled scan machinery of \cref{sec:results-scan}.

\begin{notebox}
\textbf{Known limitations, stated explicitly.} The systematic covariance
omits NC's own two dedicated systematics (Table XIX/XX, never
transcribed) and two Table XXXIV rows' CC$_9$-CC$_{17}$ entries
(illegible in the rendered source PDF, set to zero and flagged rather than
guessed, \cref{ch:data-release}). No combination with KamLAND or other
solar experiments is attempted, so $\Delta m^2_{21}$ remains only weakly
constrained here, matching SNO's own quoted SNO-only precision (Table
XXVIII). Next steps: verify the flagged Table XXXIV entries against the
primary source directly, transcribe Table XIX/XX's NC-specific
systematics, and combine with this project's own real Borexino/Daya Bay
fits for a joint solar-sector constraint.
\end{notebox}
